
\documentclass{article}
\usepackage{colortbl}
\usepackage{makecell}
\usepackage{multirow}
\usepackage{supertabular}

\begin{document}

\newcounter{utterance}

\centering \large Interaction Transcript for game `dond', experiment `coop\_en', episode 1 with claude{-}sonnet{-}4{-}5{-}azure{-}high{-}t1.0.
\vspace{24pt}

{ \footnotesize  \setcounter{utterance}{1}
\setlength{\tabcolsep}{0pt}
\begin{supertabular}{c@{$\;$}|p{.15\linewidth}@{}p{.15\linewidth}p{.15\linewidth}p{.15\linewidth}p{.15\linewidth}p{.15\linewidth}}
    \# & $\;$A & \multicolumn{4}{c}{Game Master} & $\;\:$B\\
    \hline

    \theutterance \stepcounter{utterance}  
    & & \multicolumn{4}{p{0.6\linewidth}}{
        \cellcolor[rgb]{0.9,0.9,0.9}{
            \makecell[{{p{\linewidth}}}]{
                \texttt{\tiny{[P1$\langle$GM]}}
                \texttt{You are playing a cooperative negotiation game in which you have to agree on how to divide a set of items among you and another player.} \\
\\ 
\texttt{Rules:} \\
\texttt{(a) You and the other player are given a set of items. Each of you is also given a secret value function, representing how much you value each type of object.} \\
\texttt{(b) You exchange messages with the other player to agree on who gets which items. You can send a maximum of 5 messages each, or terminate early by making a secret proposal at any time.} \\
\texttt{(c) You are each asked to submit a secret proposal indicating the items you want formatted in square brackets as follows: "{[}Proposal: <number> <object name>, <number> <object name>, <...>{]}"} \\
\texttt{(d) If your proposals are complementary, i.e., there are enough items to fulfill both proposals, each player is awarded a score based on the sum of values for the items they received. Otherwise, both of you get zero points.} \\
\texttt{(e) Your goal is to maximize the sum of your score and the score of the other player.} \\
\\ 
\texttt{Let us start.} \\
\\ 
\texttt{The set of available items is:} \\
\\ 
\texttt{2 bells, 1 doll, 2 threads} \\
\\ 
\texttt{Your secret value function is:} \\
\\ 
\texttt{bell: 1, doll: 2, thread: 3} \\
\\ 
\texttt{IMPORTANT: Your messages, unless it is the secret proposal, are directly transmitted to the other player, so do not include any response to the rules or text announcing your message. To make a secret proposal, use the indicated format. Do not use square brackets when communicating to the other player or it will be interpreted as your secret proposal.} \\
            }
        }
    }
    & \\ \\

    \theutterance \stepcounter{utterance}  
    & \multicolumn{4}{p{0.6\linewidth}}{
        \cellcolor[rgb]{0.9,0.9,0.9}{
            \makecell[{{p{\linewidth}}}]{
                \texttt{\tiny{[P1$\rangle$GM]}}
                \texttt{Hello! I'd like to work together to find the best division of items for both of us. To maximize our combined score, we should each focus on the items we value most. Could you share which items are most valuable to you?} \\
            }
        }
    }
    & & \\ \\

    \theutterance \stepcounter{utterance}  
    & & & \multicolumn{2}{p{0.3\linewidth}}{
        \cellcolor[rgb]{0.9,0.9,0.9}{
            \makecell[{{p{\linewidth}}}]{
                \texttt{\tiny{[GM$|$GM]}}
                \texttt{continue} \\
            }
        }
    }
    & & \\ \\

    \theutterance \stepcounter{utterance}  
    & & & \multicolumn{2}{p{0.3\linewidth}}{
        \cellcolor[rgb]{0.9,0.9,0.9}{
            \makecell[{{p{\linewidth}}}]{
                \texttt{\tiny{[GM$|$GM]}}
                \texttt{Hello! I'd like to work together to find the best division of items for both of us. To maximize our combined score, we should each focus on the items we value most. Could you share which items are most valuable to you?} \\
            }
        }
    }
    & & \\ \\

    \theutterance \stepcounter{utterance}  
    & & \multicolumn{4}{p{0.6\linewidth}}{
        \cellcolor[rgb]{0.9,0.9,0.9}{
            \makecell[{{p{\linewidth}}}]{
                \texttt{\tiny{[GM$\rangle$P2]}}
                \texttt{You are playing a cooperative negotiation game in which you have to agree on how to divide a set of items among you and another player.} \\
\\ 
\texttt{Rules:} \\
\texttt{(a) You and the other player are given a set of items. Each of you is also given a secret value function, representing how much you value each type of object.} \\
\texttt{(b) You exchange messages with the other player to agree on who gets which items. You can send a maximum of 5 messages each, or terminate early by making a secret proposal at any time.} \\
\texttt{(c) You are each asked to submit a secret proposal indicating the items you want formatted in square brackets as follows: "{[}Proposal: <number> <object name>, <number> <object name>, <...>{]}"} \\
\texttt{(d) If your proposals are complementary, i.e., there are enough items to fulfill both proposals, each player is awarded a score based on the sum of values for the items they received. Otherwise, both of you get zero points.} \\
\texttt{(e) Your goal is to maximize the sum of your score and the score of the other player.} \\
\\ 
\texttt{Let us start.} \\
\\ 
\texttt{The set of available items is:} \\
\\ 
\texttt{2 bells, 1 doll, 2 threads} \\
\\ 
\texttt{Your secret value function is:} \\
\\ 
\texttt{bell: 0, doll: 2, thread: 4} \\
\\ 
\texttt{IMPORTANT: Your messages, unless it is the secret proposal, are directly transmitted to the other player, so do not include any response to the rules or text announcing your message. To make a secret proposal, use the indicated format. Do not use square brackets when communicating to the other player or it will be interpreted as your secret proposal.} \\
\\ 
\\ 
\texttt{Hello! I'd like to work together to find the best division of items for both of us. To maximize our combined score, we should each focus on the items we value most. Could you share which items are most valuable to you?} \\
            }
        }
    }
    & \\ \\

    \theutterance \stepcounter{utterance}  
    & & & \multicolumn{4}{p{0.6\linewidth}}{
        \cellcolor[rgb]{0.9,0.9,0.9}{
            \makecell[{{p{\linewidth}}}]{
                \texttt{\tiny{[GM$\langle$P2]}}
                \texttt{I'm happy to collaborate! For me, threads are the most valuable items, followed by the doll. I don't value bells at all, so if bells are useful to you, you're welcome to both of them. What about you {-} which items do you value most?} \\
            }
        }
    }
     \\ \\

    \theutterance \stepcounter{utterance}  
    & & & \multicolumn{2}{p{0.3\linewidth}}{
        \cellcolor[rgb]{0.9,0.9,0.9}{
            \makecell[{{p{\linewidth}}}]{
                \texttt{\tiny{[GM$|$GM]}}
                \texttt{continue} \\
            }
        }
    }
    & & \\ \\

    \theutterance \stepcounter{utterance}  
    & & & \multicolumn{2}{p{0.3\linewidth}}{
        \cellcolor[rgb]{0.9,0.9,0.9}{
            \makecell[{{p{\linewidth}}}]{
                \texttt{\tiny{[GM$|$GM]}}
                \texttt{I'm happy to collaborate! For me, threads are the most valuable items, followed by the doll. I don't value bells at all, so if bells are useful to you, you're welcome to both of them. What about you {-} which items do you value most?} \\
            }
        }
    }
    & & \\ \\

    \theutterance \stepcounter{utterance}  
    & & \multicolumn{4}{p{0.6\linewidth}}{
        \cellcolor[rgb]{0.9,0.9,0.9}{
            \makecell[{{p{\linewidth}}}]{
                \texttt{\tiny{[P1$\langle$GM]}}
                \texttt{I'm happy to collaborate! For me, threads are the most valuable items, followed by the doll. I don't value bells at all, so if bells are useful to you, you're welcome to both of them. What about you {-} which items do you value most?} \\
            }
        }
    }
    & \\ \\

    \theutterance \stepcounter{utterance}  
    & \multicolumn{4}{p{0.6\linewidth}}{
        \cellcolor[rgb]{0.9,0.9,0.9}{
            \makecell[{{p{\linewidth}}}]{
                \texttt{\tiny{[P1$\rangle$GM]}}
                \texttt{Thank you for sharing! Interestingly, I also value threads the most at 3 points each, followed by the doll at 2 points, and bells at 1 point each. } \\
\\ 
\texttt{Since we both value threads highly, we need to figure out the optimal split. Given that you don't value bells at all and I do, I should definitely take both bells (worth 2 points to me, 0 to you).} \\
\\ 
\texttt{For the threads and doll {-} since we both value threads at our highest level, could you tell me your exact point values? That way we can calculate which division gives us the maximum combined score. For example, if your thread value is higher than 3, you should get both threads. If it's lower, I should get them. And we can figure out the doll similarly.} \\
            }
        }
    }
    & & \\ \\

    \theutterance \stepcounter{utterance}  
    & & & \multicolumn{2}{p{0.3\linewidth}}{
        \cellcolor[rgb]{0.9,0.9,0.9}{
            \makecell[{{p{\linewidth}}}]{
                \texttt{\tiny{[GM$|$GM]}}
                \texttt{continue} \\
            }
        }
    }
    & & \\ \\

    \theutterance \stepcounter{utterance}  
    & & & \multicolumn{2}{p{0.3\linewidth}}{
        \cellcolor[rgb]{0.9,0.9,0.9}{
            \makecell[{{p{\linewidth}}}]{
                \texttt{\tiny{[GM$|$GM]}}
                \texttt{Thank you for sharing! Interestingly, I also value threads the most at 3 points each, followed by the doll at 2 points, and bells at 1 point each. } \\
\\ 
\texttt{Since we both value threads highly, we need to figure out the optimal split. Given that you don't value bells at all and I do, I should definitely take both bells (worth 2 points to me, 0 to you).} \\
\\ 
\texttt{For the threads and doll {-} since we both value threads at our highest level, could you tell me your exact point values? That way we can calculate which division gives us the maximum combined score. For example, if your thread value is higher than 3, you should get both threads. If it's lower, I should get them. And we can figure out the doll similarly.} \\
            }
        }
    }
    & & \\ \\

    \theutterance \stepcounter{utterance}  
    & & \multicolumn{4}{p{0.6\linewidth}}{
        \cellcolor[rgb]{0.9,0.9,0.9}{
            \makecell[{{p{\linewidth}}}]{
                \texttt{\tiny{[GM$\rangle$P2]}}
                \texttt{Thank you for sharing! Interestingly, I also value threads the most at 3 points each, followed by the doll at 2 points, and bells at 1 point each. } \\
\\ 
\texttt{Since we both value threads highly, we need to figure out the optimal split. Given that you don't value bells at all and I do, I should definitely take both bells (worth 2 points to me, 0 to you).} \\
\\ 
\texttt{For the threads and doll {-} since we both value threads at our highest level, could you tell me your exact point values? That way we can calculate which division gives us the maximum combined score. For example, if your thread value is higher than 3, you should get both threads. If it's lower, I should get them. And we can figure out the doll similarly.} \\
            }
        }
    }
    & \\ \\

    \theutterance \stepcounter{utterance}  
    & & & \multicolumn{4}{p{0.6\linewidth}}{
        \cellcolor[rgb]{0.9,0.9,0.9}{
            \makecell[{{p{\linewidth}}}]{
                \texttt{\tiny{[GM$\langle$P2]}}
                \texttt{You're absolutely right to think about optimization! Here are my exact values:} \\
\\ 
\texttt{{-} Threads: 4 points each} \\
\texttt{{-} Doll: 2 points} \\
\texttt{{-} Bells: 0 points} \\
\\ 
\texttt{So since I value threads at 4 points and you value them at 3 points, I should get both threads to maximize our combined score. You should definitely get both bells (worth 2 to you, 0 to me).} \\
\\ 
\texttt{For the doll, we both value it at 2 points, so it doesn't matter for the total which of us takes it. The maximum combined score is 12 points, achieved by:} \\
\texttt{{-} Either: I take both threads + doll (10 points), you take both bells (2 points)} \\
\texttt{{-} Or: I take both threads (8 points), you take both bells + doll (4 points)} \\
\\ 
\texttt{Which would you prefer? Both give us the same total of 12 points.} \\
            }
        }
    }
     \\ \\

    \theutterance \stepcounter{utterance}  
    & & & \multicolumn{2}{p{0.3\linewidth}}{
        \cellcolor[rgb]{0.9,0.9,0.9}{
            \makecell[{{p{\linewidth}}}]{
                \texttt{\tiny{[GM$|$GM]}}
                \texttt{continue} \\
            }
        }
    }
    & & \\ \\

    \theutterance \stepcounter{utterance}  
    & & & \multicolumn{2}{p{0.3\linewidth}}{
        \cellcolor[rgb]{0.9,0.9,0.9}{
            \makecell[{{p{\linewidth}}}]{
                \texttt{\tiny{[GM$|$GM]}}
                \texttt{You're absolutely right to think about optimization! Here are my exact values:} \\
\\ 
\texttt{{-} Threads: 4 points each} \\
\texttt{{-} Doll: 2 points} \\
\texttt{{-} Bells: 0 points} \\
\\ 
\texttt{So since I value threads at 4 points and you value them at 3 points, I should get both threads to maximize our combined score. You should definitely get both bells (worth 2 to you, 0 to me).} \\
\\ 
\texttt{For the doll, we both value it at 2 points, so it doesn't matter for the total which of us takes it. The maximum combined score is 12 points, achieved by:} \\
\texttt{{-} Either: I take both threads + doll (10 points), you take both bells (2 points)} \\
\texttt{{-} Or: I take both threads (8 points), you take both bells + doll (4 points)} \\
\\ 
\texttt{Which would you prefer? Both give us the same total of 12 points.} \\
            }
        }
    }
    & & \\ \\

    \theutterance \stepcounter{utterance}  
    & & \multicolumn{4}{p{0.6\linewidth}}{
        \cellcolor[rgb]{0.9,0.9,0.9}{
            \makecell[{{p{\linewidth}}}]{
                \texttt{\tiny{[P1$\langle$GM]}}
                \texttt{You're absolutely right to think about optimization! Here are my exact values:} \\
\\ 
\texttt{{-} Threads: 4 points each} \\
\texttt{{-} Doll: 2 points} \\
\texttt{{-} Bells: 0 points} \\
\\ 
\texttt{So since I value threads at 4 points and you value them at 3 points, I should get both threads to maximize our combined score. You should definitely get both bells (worth 2 to you, 0 to me).} \\
\\ 
\texttt{For the doll, we both value it at 2 points, so it doesn't matter for the total which of us takes it. The maximum combined score is 12 points, achieved by:} \\
\texttt{{-} Either: I take both threads + doll (10 points), you take both bells (2 points)} \\
\texttt{{-} Or: I take both threads (8 points), you take both bells + doll (4 points)} \\
\\ 
\texttt{Which would you prefer? Both give us the same total of 12 points.} \\
            }
        }
    }
    & \\ \\

    \theutterance \stepcounter{utterance}  
    & \multicolumn{4}{p{0.6\linewidth}}{
        \cellcolor[rgb]{0.9,0.9,0.9}{
            \makecell[{{p{\linewidth}}}]{
                \texttt{\tiny{[P1$\rangle$GM]}}
                \texttt{Perfect analysis! Both options do give us the maximum combined score of 12 points. I'd prefer the second option where I take both bells and the doll (4 points for me) and you take both threads (8 points for you). This gives me a slightly better individual outcome while maintaining our optimal total.} \\
\\ 
\texttt{Shall we go with that division?} \\
\\ 
\texttt{{[}Proposal: 2 bells, 1 doll{]}} \\
            }
        }
    }
    & & \\ \\

    \theutterance \stepcounter{utterance}  
    & & & \multicolumn{2}{p{0.3\linewidth}}{
        \cellcolor[rgb]{0.9,0.9,0.9}{
            \makecell[{{p{\linewidth}}}]{
                \texttt{\tiny{[GM$|$GM]}}
                \texttt{proposal} \\
            }
        }
    }
    & & \\ \\

    \theutterance \stepcounter{utterance}  
    & & & \multicolumn{2}{p{0.3\linewidth}}{
        \cellcolor[rgb]{0.9,0.9,0.9}{
            \makecell[{{p{\linewidth}}}]{
                \texttt{\tiny{[GM$|$GM]}}
                [2, 1, 0]
            }
        }
    }
    & & \\ \\

    \theutterance \stepcounter{utterance}  
    & & \multicolumn{4}{p{0.6\linewidth}}{
        \cellcolor[rgb]{0.9,0.9,0.9}{
            \makecell[{{p{\linewidth}}}]{
                \texttt{\tiny{[GM$\rangle$P2]}}
                \texttt{The other player has submitted a secret proposal. It is now your turn to submit a proposal of your own.} \\
            }
        }
    }
    & \\ \\

    \theutterance \stepcounter{utterance}  
    & & & \multicolumn{4}{p{0.6\linewidth}}{
        \cellcolor[rgb]{0.9,0.9,0.9}{
            \makecell[{{p{\linewidth}}}]{
                \texttt{\tiny{[GM$\langle$P2]}}
                \texttt{{[}Proposal: 2 threads{]}} \\
            }
        }
    }
     \\ \\

    \theutterance \stepcounter{utterance}  
    & & & \multicolumn{2}{p{0.3\linewidth}}{
        \cellcolor[rgb]{0.9,0.9,0.9}{
            \makecell[{{p{\linewidth}}}]{
                \texttt{\tiny{[GM$|$GM]}}
                \texttt{proposal} \\
            }
        }
    }
    & & \\ \\

    \theutterance \stepcounter{utterance}  
    & & & \multicolumn{2}{p{0.3\linewidth}}{
        \cellcolor[rgb]{0.9,0.9,0.9}{
            \makecell[{{p{\linewidth}}}]{
                \texttt{\tiny{[GM$|$GM]}}
                [0, 0, 2]
            }
        }
    }
    & & \\ \\

    \theutterance \stepcounter{utterance}  
    & & & \multicolumn{2}{p{0.3\linewidth}}{
        \cellcolor[rgb]{0.9,0.9,0.9}{
            \makecell[{{p{\linewidth}}}]{
                \texttt{\tiny{[GM$|$GM]}}
                [[2, 1, 0], [0, 0, 2]]
            }
        }
    }
    & & \\ \\

\end{supertabular}
}

\end{document}
